% sections/12_conclusion.tex -- Section 12, Conclusion (W4). Plan: BLUEPRINT 2.12; one page; no equations; no new numbers.
% The only digits are the saving-subsidy lower bounds already printed in Table 9 and Section 9.3 (cell cited below).
\section{Conclusion}\label{sec:conclusion}

Social Security's financing gap can be addressed through the claiming-age schedule in two ways, and the estimated life-cycle model of this paper assigns them different welfare signs. A higher normal retirement age lowers the benefit at every claiming age; it reduces the cohort's benefit outlays and net government spending before anyone changes behavior, and it lowers ex ante welfare, so the welfare delivered per dollar of net spending by reversing it is small. A steeper reward for delayed claiming, through a higher delayed retirement credit or a symmetric adjustment around the normal retirement age, leaves the benefit at the normal retirement age unchanged and changes the price of moving the claiming age. In the model it raises ex ante welfare and lowers both benefit outlays and net government spending, so its marginal value of public funds is infinite, and it moves participation from the years before the normal retirement age to the years after it. At approximately equal fiscal savings the two reforms differ in the sign of their welfare effect. The steepest adjustments shift the burden toward individuals who claim early, and the effects on them are not described by the average welfare measure.

The same accounting ranks the other instruments. Among increases in the generosity of retirement benefits, raising the first bracket of the benefit formula delivers the most welfare per dollar of net spending, and the first step is approximately self-financing; proportional increases deliver less, and increases concentrated in the upper brackets the least, while lowering participation at older ages the most. Expansions of Medicare coverage after age 65 raise welfare and lower participation at ages 62--69, reductions in income taxes after age 62 raise welfare and participation at every age, and both return little per dollar of net spending. Subsidies to retirement saving return at least 0.84--0.96 dollars of welfare per dollar of subsidy, with an upper bound that depends on whether the induced asset income is taxed as it accrues. % src: T:welfare_savings_mvpf_v3.tex:120,122,123
% src: T:welfare_medicare_mvpfv8.tex:113,116 (62-64: -1.70, -1.95; 65-69: -0.38, -3.44; 30-49: 0.22, 0.21); T:welfare_tax_mvpf_v8.tex:114-118 (all three participation columns positive)
Across the menu, the sign of the net fiscal effect, not the size of the welfare gain, separates the reforms that contribute to closing the financing gap from those that widen it.

The exercise has defined limits: a partial-equilibrium evaluation for one cohort at fixed factor prices, with no transition and no weights on other cohorts, with statutory ages acting only through incentives, an individual decision-maker, accidental bequests, and one parameter vector and one preference specification; \autoref{sec:scope} states the direction in which each bears on the results, and the marginal value of public funds is reported so that other weights on public funds can be applied. This paper evaluates reforms within the rules as they are written; the optimum against which those rules could be judged is characterized in the companion paper \citep{ndiaye2026flexible}.
